\documentclass[12pt]{article}

\usepackage[a4paper, margin=2.5cm]{geometry}
\usepackage[utf8]{inputenc}                             % UTF8 enkodiranje
\usepackage[slovene]{babel}                             % Slovenščina
\usepackage[pdfusetitle, hidelinks, unicode]{hyperref}  % Nastavi atribute PDF-ja, ne označuj povezav
\usepackage{microtype}                                  % Podzavestne izboljšave za tipografijsko perfekcijo :)
\usepackage{enumitem}                                   % Seznami za člene
\usepackage{graphicx}                                   % Vključitev slik
\usepackage{dirtytalk}                                  % Citat
\usepackage{listings}                                   % Kodni blok
\usepackage{fancyvrb}
\usepackage[font=]{caption}                             % Required for specifying captions
\usepackage[normalem]{ulem}                             % Krašanje enot v enačbi
\usepackage{times}                                      % Times New Roman pisava
\usepackage{tikz} 
\usepackage[european]{circuitikz}                       % Električna vezja
\usepackage{datetime}                                   % Datum
\usepackage{csquotes}
\usepackage[style=ieee, maxbibnames=3, minbibnames=1, maxcitenames=1, mincitenames=1]{biblatex}                       % Navajanje virov
\addbibresource{viri.bib}

\setlength{\parindent}{0em}
\setlength{\parskip}{1ex}

\setcounter{secnumdepth}{5}
\setcounter{tocdepth}{5}

\renewcommand{\thesection}{\arabic{section}}
\renewcommand{\thesubsection}{\thesection.\arabic{subsection}}
\renewcommand{\thesubsubsection}{\thesubsection.\arabic{subsubsection}}
\renewcommand{\theparagraph}{\thesubsubsection.\arabic{paragraph}}
\renewcommand{\thesubparagraph}{\theparagraph.\arabic{subparagraph}}

\renewcommand{\labelnamepunct}{\addcomma\space}
\DeclareFieldFormat[article]{title}{#1}
\DeclareFieldFormat[online]{title}{\mkbibemph{#1}}

\DefineBibliographyStrings{slovene}{
  andothers = {et. al\adddot},
  urlseen = {dostopano:}
}

\newdateformat{MMYYYYdate}{\monthname[\THEMONTH] \THEYEAR}

\author{Jaka Kovač, G 3. b}

\begin{document}
\pagenumbering{arabic}

\begin{center}
    \thispagestyle{empty}
    \includegraphics[scale=1]{slike/logotip_vegova_leze_brezokvirja.png}
    
    \vspace{\fill} 
    Strokovno poročilo pri predmetu elektrotehnika

    \title{Komunikacijska vezja in naprave}

    \normalsize
    Podnaslov seminarske naloge k si ga še nism izmislu
    \vspace{\fill}

    Mentor: Anton Orehek, uni. dipl. inž., prof. \hfill Avtor: Jaka Kovač, G 3. b\\
    \null
    Ljubljana, december 2022 - \MMYYYYdate\today
\end{center}
\newpage
\null
\newpage

%If I had more time I would have written a shorter letter - Mark Twain
\section*{Povzetek}
\section*{Abstract}

% KAZALO 
\newpage
\tableofcontents

\newpage
\begingroup
\makeatletter
\section*{Slike}
\@starttoc{lof}
\let\clearpage\relax
\section*{Tabele}
\@starttoc{lot}
\makeatother
\endgroup


\newpage
\section{Analogna komunikacija}
    \subsection{Začetki radia (in radioamaterstva)}
        \subsubsection{Zgodovina radia}
            Prav vsi poznamo radio. To je tista majhna naprava v avtu, ki voznikom (in potnikom) olajša čas, ki so ga prisiljeni preživeti za volanom. 
            Veliko ljudi pa se ne zaveda, da je radio mnogo več. Slovar slovenskega knjižnega jezika s prvim pomenom definira radio kot \textit{naprava za oddajanje in sprejemanje električnih impulzov, signalov po radijskih valovih}. \cite{SSKJ-radio}\\
            Leta 1895 \cite{ppt} je potekal prvi prenos sporočila s pomočjo radijskih valov, osem let kasneje pa prva uspešna (enosmerna) komunikacija iz ZDA v Združeno kraljestvo. 
            Leta 1920 sta v ZDA in Veliki Britaniji pričeli delovati prvi radiodifuzni\footnote{radiodifuzija – oddajanje radijskih signalov namenjenih poslušanju} postaji, leta 1928 pa je Radio Ljubljana postala prva radiodifuzna postaja v Sloveniji.
        \subsubsection{Radioamaterstvo}
            V \textit{Priročniku za radioamaterje} je RADIOAMATERSTVO predstavljeno z: \textit{"}Zelo poenostavljeno bi lahko rekli, da je to ljubiteljsko, nepoklicno ukvarjanje z radiom oziroma radiotehniko." \cite{HAM-prirocnik}
            Ker sem tudi sam radioamater vem, da ta dejavnost pomeni mnogo več. \\
            Za legalno delovanje morajo radioamaterji opraviti izpit, ena izmed dejavnosti aktivnih radioamaterjev je zato izobraževanje.
            Druga dejavnost so prirejanje in udeležba tekmovanj. Zveza radioamaterjev Slovenije podeljuje dve t.i. diplomi.\\
            Radioamaterji po potrebi tudi pomagajo pri večjih nesrečah. To jim narekuje kodeks ARON (aktivnosti radioamaterjev ob nesrečah in nevarnostih), ki se ga je nazadnje množično aktiviralo ob potresu leta 2020 (pred tem pa ob žledu)
    \subsection{Analogni signali}
            Analogni signali so tisti signali, ki lahko zavzamejo vse vrednosti na določenem intervalu. Čas je primer analogne vrednosti, ker mu ne moremo odločiti najmanjše enote za katero bi se spremenil. Urni kazalec se premika s stalno hitrostjo. To pomeni da se v neskončno majhnem intervalu časa vseeno spremeni za nek delež stopinje, vendar pa ljudje tega navadno ne opazimo. 
        \subsubsection{Amplitudna modulacija}
            \paragraph{Nek po poglavje}
            \paragraph{Nek drug po poglavje}
        \subsubsection{Frekvenčna modulacija}
            \paragraph{Nek po poglavje}
            \paragraph{Nek drug po poglavje}

    \newpage
    \subsection{Fourierova transformacija}
        \subsubsection{FFT – hitra fourierova transformacija}
    \subsection{Problemi analognih komunikacij}
        Tole je pomembno besedilo

\begin{table}[h]
    \centering
    \caption{This is a table}
    \begin{tabular}{lcr}
      \hline
      Column 1 & Column 2 & Column 3 \\
      \hline
      A & B & C \\
      D & E & F \\
      \hline
    \end{tabular}
    \label{tab:table}
\end{table}  


\newpage
\section{Digitalna komunikacija}
    \subsection{Prednosti in slabosti digitalne komunikacije}
    \subsection{Problemi digitalnih komunikacij}
        \subsubsection{Bitflips}
    \subsection{Rešitve problemov}
        \subsubsection{Error correction}
            \paragraph{Reed-Solomon kod}
            \paragraph{tisto s polinomi in črno magijo}

\newpage
\section{Empirični del}
    \subsection{Analogna vezja}


\newpage
    \subsection{Digitalna vezja}
        \begin{figure}[h!]
            \begin{center}
                \caption{Moje prvo \LaTeX vezje}
                \begin{circuitikz} \draw
                    (0, 2.5) to[american voltage source, l_=$U_{vh}$, i<=$I_{vh}$] (0, 0)
                    (0, 2.5) to[D, l=$D_1$] (3.5, 2.5)
                    to[R, l=$R_1$] (3.5, 0)
                    to[short] (0, 0)
                    ;
                \end{circuitikz}
                \label{fig:vezje1}
            \end{center}
        \end{figure}

\newpage

\begingroup
    \makeatletter
        \section{Viri in literatura}
            \nocite{*}
            \printbibliography[heading=none]
    \makeatother
\endgroup
\end{document}